% ============================================================
%  Chapter 10: Thermodynamics
% ============================================================
\chapter{热力学：从统计到宏观}
\label{ch:thermodynamics}

\section{概述}

R6（$S = k_B \ln W$）+ 等概率先验假设构成了统计力学的全部基础。本章从 R6 出发推导热力学第零/一/三定律、理想气体状态方程、Fermi-Dirac 和 Bose-Einstein 量子统计分布。

% ════════════════════════════════════════════════════════
\section{热力学第零定律}
% ════════════════════════════════════════════════════════

\begin{lawbox}[热力学第零定律]
  $T_A = T_B,\; T_B = T_C \implies T_A = T_C$（温度传递性）
  \quad\textbf{前提}：R6（$S = k_B\ln W$, $\Delta S \ge 0$）
\end{lawbox}

\begin{derivationbox}[推导]
  \textbf{Step 1} — 两系统热接触达到平衡时熵最大化：
  $dS_{\text{total}} = dS_A + dS_B = (1/T_A)dE_A + (1/T_B)dE_B = 0$
  \hfill\textit{(R6: $dS = dE/T$)}

  \textbf{Step 2} — 能量守恒 $dE_A + dE_B = 0$：
  $(1/T_A - 1/T_B)dE_A = 0 \implies T_A = T_B$
  \hfill\texttt{[N]} 孤立系统（$dE_{\text{total}}=0$）

  \textbf{Step 3} — 传递性：$T_A = T_B$ 且 $T_B = T_C \implies T_A = T_C$（等式传递性）
\end{derivationbox}

% ════════════════════════════════════════════════════════
\section{热力学第一定律}
% ════════════════════════════════════════════════════════

\begin{lawbox}[热力学第一定律]
  $dU = T\,dS - P\,dV + \mu\,dN$
  \quad\textbf{前提}：能量守恒（Noether）, R6
\end{lawbox}

\noindent\textbf{变量}：内能 $U$ [M$\cdot$L$^2\cdot$T$^{-2}$]（状态函数），压强 $P$ [M$\cdot$L$^{-1}\cdot$T$^{-2}$],
热 $Q$ [M$\cdot$L$^2\cdot$T$^{-2}$]（过程量）, 功 $W$ [M$\cdot$L$^2\cdot$T$^{-2}$]（过程量）。

\begin{derivationbox}[推导]
  \textbf{Step 1} — 能量守恒：$dU = \delta Q - \delta W$
  \hfill\textit{（$dU$ 全微分；$\delta Q,\delta W$ 非全微分）}

  \textbf{Step 2} — 可逆过程 $\delta Q_{\text{rev}} = T\,dS$
  \hfill\texttt{[N]} 可逆过程（R6: $dS = \delta Q_{\text{rev}}/T$）

  \textbf{Step 3} — 简单流体 $\delta W = P\,dV$

  \textbf{Step 4} — 推广到可变粒子数：$dU = T\,dS - P\,dV + \mu\,dN$
\end{derivationbox}

% ════════════════════════════════════════════════════════
\section{热力学第三定律}
% ════════════════════════════════════════════════════════

\begin{lawbox}[热力学第三定律]
  $S \to 0$ as $T \to 0$（完美晶体，非简并基态）
  \quad\textbf{前提}：R6, 量子离散能谱（$\leftarrow$ Schr\"{o}dinger 方程）
\end{lawbox}

\noindent\textbf{变量}：基态简并度 $g_0$ [1]（能量最低能级的独立量子态数目）。

\begin{derivationbox}[推导]
  \textbf{Step 1} — $S = k_B \ln W$ (R6)

  \textbf{Step 2} — $T \to 0$：系统进入量子基态。
  离散能谱由 Schr\"{o}dinger 方程给出（$\leftarrow$ R4b+R4c）
  \hfill\texttt{[N]} 量子力学（离散能谱）——经典系统 $S \to -\infty$（Gibbs 悖论）

  \textbf{Step 3} — 若基态非简并（$g_0=1$）：$W=1 \implies S = k_B\ln 1 = 0$
  若基态 $g_0$ 重简并：$S \to k_B\ln g_0 > 0$（剩余熵，如水冰 Ih）
\end{derivationbox}

% ════════════════════════════════════════════════════════
\section{理想气体状态方程}
% ════════════════════════════════════════════════════════

\begin{lawbox}[理想气体状态方程]
  $PV = nRT$（$R = N_A k_B = 8.314$ J$\cdot$mol$^{-1}\cdot$K$^{-1}$）
  \quad\textbf{前提}：R6, 热力学第一定律
\end{lawbox}

\begin{derivationbox}[推导]
  \textbf{Step 1} — 自由粒子在体积 $V$ 中：$W \propto V^N$
  \hfill\textit{（经典稀薄气体，无相互作用）}

  \textbf{Step 2} — $S = k_B\ln W = Nk_B\ln V + C(E,N)$

  \textbf{Step 3} — 从热力学基本方程求压强：
  $P = T(\partial S/\partial V)_{E,N} = T\cdot(Nk_B/V)$

  \textbf{Step 4} — $PV = Nk_B T = nRT\qquad$\texttt{[N]} R1（$D=3 \implies V \propto L^3$）
\end{derivationbox}

% ════════════════════════════════════════════════════════
\section{Fermi-Dirac 分布（推论）}
% ════════════════════════════════════════════════════════

\begin{lawbox}[Fermi-Dirac 分布 — corollary]
  $f(E) = \dfrac{1}{\exp\!\big(\frac{E-\mu}{k_B T}\big) + 1}$
  \quad\textbf{前提}：Pauli 不相容 ($n_i \in \{0,1\}$), R6
\end{lawbox}

\begin{derivationbox}[推导]
  \textbf{Step 1} — Fermi 子约束：$n_i \in \{0,1\}$（Pauli 不相容原理）

  \textbf{Step 2} — 最大化熵 $S = -k_B\sum_i[n_i\ln n_i + (1-n_i)\ln(1-n_i)]$\\
  约束 $\sum n_i = N$，$\sum n_i E_i = E$

  \textbf{Step 3} — Lagrange 乘子法：$\partial/\partial n_i[\cdots] = 0$
  $\implies \ln\frac{1-n_i}{n_i} = \alpha + \beta E_i$

  \textbf{Step 4} — $\alpha = -\beta\mu$，$\beta = 1/(k_B T)$：
  $n_i = \dfrac{1}{\exp((E_i-\mu)/k_B T) + 1}$
\end{derivationbox}

% ════════════════════════════════════════════════════════
\section{Bose-Einstein 分布（推论）}
% ════════════════════════════════════════════════════════

\begin{lawbox}[Bose-Einstein 分布 — corollary]
  $f(E) = \dfrac{1}{\exp\!\big(\frac{E-\mu}{k_B T}\big) - 1}$
  \quad\textbf{前提}：叠加原理 ($n_i \in \{0,1,2,\dots\}$), R6
\end{lawbox}

\begin{derivationbox}[推导]
  \textbf{Step 1} — Bose 子约束：$n_i \in \{0,1,2,\dots\}$（无上限——R4a 叠加原理）

  \textbf{Step 2} — 最大化熵 $S = k_B\sum_i[(1+n_i)\ln(1+n_i) - n_i\ln n_i]$

  \textbf{Step 3} — Lagrange 乘子法 $\implies n_i = \dfrac{1}{\exp((E_i-\mu)/k_B T) - 1}$

  分母中的 $-1$（vs FD 的 $+1$）$\implies$ 当 $\mu \to 0$ 时 $f$ 发散 $\implies$ Bose-Einstein 凝聚
\end{derivationbox}
